\section{联合RMSD低于5~\AA{}的101条候选清单}
\label{app:pass-list}

表~\ref{tab:pass101}逐条列出六个非3AV复合物中同时满足全复合物CA RMSD
$<5~\angstrom$且最佳正向循环肽CA RMSD $<5~\angstrom$的101条候选。
这些候选均来自正文定义的476条保留设计，并沿用“一次全复合物叠合、完整末链
CA配对、最佳正向循环移位、不反向、不进行肽链二次拟合”的联合RMSD口径。
“类别”列进一步标出其中16条联合$<3~\angstrom$候选；其余85条位于
$[3,5)~\angstrom$类别。序列大小写保持原始输出，小写残基表示严格$p>0.6$
规则保存的预测甲基化标记，甲基化位点使用1基编号。

可复核的机器可读文件同时随稿提供：本表101条的完整记录位于
\path{data/Pass_LT5_101.csv}，其中16条联合$<3~\angstrom$候选另存于
\path{data/Pass_LT3_16.csv}；包含通过与未通过候选、各审计门控字段及两种
RMSD的完整476条分析表位于\path{data/Candidates_476.csv}。因此，本附录是便于
人工阅读的通过清单，完整统计复算应以476条CSV为准。

\begin{landscape}
\begingroup
\small
\setlength{\tabcolsep}{4pt}
\setlength\LTleft{0pt}
\setlength\LTright{0pt}
\captionof{table}{六个非 3AV 复合物中联合 RMSD 低于 5~\AA{} 的 101 条候选}
\label{tab:pass101}
\begin{longtable}{C{10mm} C{18mm} L{48mm} L{36mm} C{27mm} C{30mm} C{24mm}}
  \toprule
  编号 & 靶点 & 设计序列（小写为甲基化） & 甲基化位点 &
  \makecell{全复合物CA\\RMSD（\AA）} &
  \makecell{循环肽CA\\RMSD（\AA）} & 类别 \\
  \midrule
  \endfirsthead

  \multicolumn{7}{l}{\footnotesize 表~\ref{tab:pass101}（续）}\\
  \toprule
  编号 & 靶点 & 设计序列（小写为甲基化） & 甲基化位点 &
  \makecell{全复合物CA\\RMSD（\AA）} &
  \makecell{循环肽CA\\RMSD（\AA）} & 类别 \\
  \midrule
  \endhead

  \midrule
  \multicolumn{7}{r}{\footnotesize 续下页}\\
  \endfoot

  \bottomrule
  \endlastfoot

  1 & 1SFI & \seq{RCCrGNCGQCQFGG} & 4 & 0.536 & 2.818 & $<3$ \\
2 & 1SFI & \seq{RCGrGcCrQGcNGG} & 4,6,8,11 & 0.562 & 3.693 & $3\text{--}<5$ \\
3 & 1SFI & \seq{RCGrCQNGQCrNGG} & 4,11 & 1.120 & 3.840 & $3\text{--}<5$ \\
4 & 1SFI & \seq{RrGrGrCGQCQFMQ} & 2,4,6 & 1.483 & 3.859 & $3\text{--}<5$ \\
5 & 1SFI & \seq{CCGrGrCrQGrGGR} & 4,6,8,11 & 0.572 & 3.967 & $3\text{--}<5$ \\
6 & 1SFI & \seq{RCGrQQCrKGcQMC} & 4,8,11 & 0.817 & 4.217 & $3\text{--}<5$ \\
7 & 1SFI & \seq{RCGrGQCrKGrNGC} & 4,8,11 & 0.242 & 4.252 & $3\text{--}<5$ \\
8 & 1SFI & \seq{RrrrGQCrQCNRDG} & 2,3,4,8 & 1.805 & 4.276 & $3\text{--}<5$ \\
9 & 1SFI & \seq{RQYrGKGrQGrRMG} & 4,8,11 & 0.221 & 4.414 & $3\text{--}<5$ \\
10 & 1SFI & \seq{RGGrCcCGGGNQMG} & 4,6 & 1.068 & 4.468 & $3\text{--}<5$ \\
11 & 1SFI & \seq{RCHrCNGCQCrQGC} & 4,11 & 0.576 & 4.523 & $3\text{--}<5$ \\
12 & 1SFI & \seq{RrGcCrCrKGKFDG} & 2,4,6,8 & 0.632 & 4.532 & $3\text{--}<5$ \\
13 & 1SFI & \seq{RCGrGQCrKGrQDG} & 4,8,11 & 0.563 & 4.612 & $3\text{--}<5$ \\
14 & 1SFI & \seq{RCGrGrCrQCrQMC} & 4,6,8,11 & 0.970 & 4.626 & $3\text{--}<5$ \\
15 & 1SFI & \seq{CCrrrGCrKCNNNG} & 3,4,5,8 & 0.219 & 4.753 & $3\text{--}<5$ \\
16 & 1SFI & \seq{RrGrGrCGQCcQMG} & 2,4,6,11 & 1.854 & 4.756 & $3\text{--}<5$ \\
17 & 1SFI & \seq{RCrrrNNGGCrQGG} & 3,4,5,11 & 0.467 & 4.763 & $3\text{--}<5$ \\
18 & 1SFI & \seq{CCYrCrCrKGrFDG} & 4,6,8,11 & 0.228 & 4.776 & $3\text{--}<5$ \\
19 & 1SFI & \seq{RCGQGQCrQCrRDG} & 8,11 & 0.561 & 4.918 & $3\text{--}<5$ \\
20 & 3P8F & \seq{CrrrGrQGKGrRGC} & 2,3,4,6,11 & 0.300 & 2.075 & $<3$ \\
21 & 3P8F & \seq{GrGrGcCCQGrCMR} & 2,4,6,11 & 0.995 & 2.713 & $<3$ \\
22 & 3P8F & \seq{GrrrCrCGQCrNMG} & 2,3,4,6,11 & 0.559 & 3.255 & $3\text{--}<5$ \\
23 & 3P8F & \seq{CrhQrGCRrCQCGN} & 2,3,5,9 & 0.277 & 3.523 & $3\text{--}<5$ \\
24 & 3P8F & \seq{CCGrrQCRQCrCRC} & 4,5,11 & 0.311 & 3.807 & $3\text{--}<5$ \\
25 & 3P8F & \seq{RCrrGKCGQGrGCG} & 3,4,11 & 1.063 & 3.847 & $3\text{--}<5$ \\
26 & 3P8F & \seq{RrrrrGCGQCrCGG} & 2,3,4,5,11 & 1.556 & 3.854 & $3\text{--}<5$ \\
27 & 3P8F & \seq{RCrrrrGGQCrCMN} & 3,4,5,6,11 & 0.648 & 3.862 & $3\text{--}<5$ \\
28 & 3P8F & \seq{RrKrGQCRQGcrRG} & 2,4,11,12 & 0.522 & 3.965 & $3\text{--}<5$ \\
29 & 3P8F & \seq{RrrrGNCCKCrCDG} & 2,3,4,11 & 1.599 & 4.081 & $3\text{--}<5$ \\
30 & 3P8F & \seq{RrrrGNCGKCrRRG} & 2,3,4,11 & 1.708 & 4.303 & $3\text{--}<5$ \\
31 & 3P8F & \seq{QrrrGQCGFCrQMG} & 2,3,4,11 & 1.456 & 4.356 & $3\text{--}<5$ \\
32 & 3P8F & \seq{RrhrCKCCrCrCMG} & 2,3,4,9,11 & 0.291 & 4.363 & $3\text{--}<5$ \\
33 & 3P8F & \seq{RCGrGGCRQCrRCC} & 4,11 & 1.002 & 4.449 & $3\text{--}<5$ \\
34 & 3P8F & \seq{RrrrGKCGQGrCMG} & 2,3,4,11 & 1.105 & 4.496 & $3\text{--}<5$ \\
35 & 3P8F & \seq{GrGrGNCRQCrFMG} & 2,4,11 & 0.405 & 4.668 & $3\text{--}<5$ \\
36 & 3P8F & \seq{RCrrrNCGQCrQDG} & 3,4,5,11 & 1.323 & 4.757 & $3\text{--}<5$ \\
37 & 3P8F & \seq{CrrrGQCGQCrrMG} & 2,3,4,11,12 & 1.266 & 4.791 & $3\text{--}<5$ \\
38 & 3P8F & \seq{RrrrQNCGQGrCMG} & 2,3,4,11 & 0.289 & 4.883 & $3\text{--}<5$ \\
39 & 3P8F & \seq{RCGrGGCGNCrCGG} & 4,11 & 1.048 & 4.921 & $3\text{--}<5$ \\
40 & 3WNE & \seq{GrKWNC} & 2 & 0.933 & 1.459 & $<3$ \\
41 & 3WNE & \seq{GrGKFG} & 2 & 0.904 & 2.915 & $<3$ \\
42 & 3WNE & \seq{GrGWQG} & 2 & 0.904 & 2.946 & $<3$ \\
43 & 3WNE & \seq{GrGKGG} & 2 & 0.912 & 3.407 & $3\text{--}<5$ \\
44 & 3WNE & \seq{GrNWRG} & 2 & 1.041 & 3.868 & $3\text{--}<5$ \\
45 & 3ZGC & \seq{rEGGQNR} & 1 & 0.973 & 2.076 & $<3$ \\
46 & 3ZGC & \seq{NrGGGQR} & 2 & 1.488 & 2.231 & $<3$ \\
47 & 3ZGC & \seq{rrGGLGR} & 1,2 & 0.988 & 3.103 & $3\text{--}<5$ \\
48 & 3ZGC & \seq{rrNGGGR} & 1,2 & 0.968 & 3.233 & $3\text{--}<5$ \\
49 & 3ZGC & \seq{cGGGQDR} & 1 & 1.589 & 3.516 & $3\text{--}<5$ \\
50 & 3ZGC & \seq{rGcGQGR} & 1,3 & 1.600 & 3.759 & $3\text{--}<5$ \\
51 & 3ZGC & \seq{EcGQGMR} & 2 & 1.194 & 3.769 & $3\text{--}<5$ \\
52 & 3ZGC & \seq{cCGGGNR} & 1 & 1.476 & 4.100 & $3\text{--}<5$ \\
53 & 3ZGC & \seq{rGGGNLR} & 1 & 2.226 & 4.123 & $3\text{--}<5$ \\
54 & 3ZGC & \seq{rrGGQDR} & 1,2 & 0.919 & 4.129 & $3\text{--}<5$ \\
55 & 3ZGC & \seq{rGGGQGR} & 1 & 2.535 & 4.262 & $3\text{--}<5$ \\
56 & 3ZGC & \seq{GrGGNQR} & 2 & 1.847 & 4.281 & $3\text{--}<5$ \\
57 & 3ZGC & \seq{rGGGGMR} & 1 & 2.742 & 4.367 & $3\text{--}<5$ \\
58 & 3ZGC & \seq{cGGGQQR} & 1 & 2.443 & 4.370 & $3\text{--}<5$ \\
59 & 3ZGC & \seq{rGGGQFR} & 1 & 2.337 & 4.454 & $3\text{--}<5$ \\
60 & 3ZGC & \seq{cGGGGDR} & 1 & 2.344 & 4.456 & $3\text{--}<5$ \\
61 & 3ZGC & \seq{cGGGQGR} & 1 & 2.284 & 4.469 & $3\text{--}<5$ \\
62 & 3ZGC & \seq{rGGGGDR} & 1 & 1.651 & 4.618 & $3\text{--}<5$ \\
63 & 3ZGC & \seq{cGGGGLR} & 1 & 2.967 & 4.651 & $3\text{--}<5$ \\
64 & 3ZGC & \seq{rGDGGGR} & 1 & 1.801 & 4.686 & $3\text{--}<5$ \\
65 & 3ZGC & \seq{rGGGQCR} & 1 & 2.553 & 4.750 & $3\text{--}<5$ \\
66 & 3ZGC & \seq{cGGGKGR} & 1 & 2.274 & 4.767 & $3\text{--}<5$ \\
67 & 3ZGC & \seq{rGGGQMR} & 1 & 2.744 & 4.898 & $3\text{--}<5$ \\
68 & 4K1E & \seq{rcrrrQCrQGNCDR} & 1,2,3,4,5,8 & 0.900 & 1.916 & $<3$ \\
69 & 4K1E & \seq{rcrrCrCrFGQCMR} & 1,2,3,4,6,8 & 0.637 & 2.157 & $<3$ \\
70 & 4K1E & \seq{rcCrGrQrNGERLR} & 1,2,4,6,8 & 0.821 & 2.330 & $<3$ \\
71 & 4K1E & \seq{rcCrrrQrQCQCDR} & 1,2,4,5,6,8 & 0.694 & 2.596 & $<3$ \\
72 & 4K1E & \seq{rcrrGrCrQGhCLR} & 1,2,3,4,6,8,11 & 0.680 & 2.801 & $<3$ \\
73 & 4K1E & \seq{CcCrrrRrQCrRLC} & 2,4,5,6,8,11 & 0.659 & 2.960 & $<3$ \\
74 & 4K1E & \seq{CcGrCQCrEGFRDR} & 2,4,8 & 0.920 & 3.273 & $3\text{--}<5$ \\
75 & 4K1E & \seq{rcrrrcCrFGQCDR} & 1,2,3,4,5,6,8 & 0.676 & 3.338 & $3\text{--}<5$ \\
76 & 4K1E & \seq{rcrrCrRrQGcCWR} & 1,2,3,4,6,8,11 & 0.644 & 3.368 & $3\text{--}<5$ \\
77 & 4K1E & \seq{rcQrrrGrQCQCDR} & 1,2,4,5,6,8 & 0.646 & 3.511 & $3\text{--}<5$ \\
78 & 4K1E & \seq{rrGrGrCrFCLRGR} & 1,2,4,6,8 & 1.221 & 3.661 & $3\text{--}<5$ \\
79 & 4K1E & \seq{rrrrGNQrQGQCMC} & 1,2,3,4,8 & 0.690 & 3.672 & $3\text{--}<5$ \\
80 & 4K1E & \seq{rcrrCrQrrCQCMR} & 1,2,3,4,6,8,9 & 0.665 & 3.714 & $3\text{--}<5$ \\
81 & 4K1E & \seq{rrrrCrCrRCQCMR} & 1,2,3,4,6,8 & 0.697 & 3.818 & $3\text{--}<5$ \\
82 & 4K1E & \seq{rcrrrrCrQCQRDR} & 1,2,3,4,5,6,8 & 1.432 & 3.950 & $3\text{--}<5$ \\
83 & 4K1E & \seq{rcCrrGQrrCQrDR} & 1,2,4,5,8,9,12 & 0.691 & 4.045 & $3\text{--}<5$ \\
84 & 4K1E & \seq{rcNrrrCrQGQCDR} & 1,2,4,5,6,8 & 1.627 & 4.420 & $3\text{--}<5$ \\
85 & 4K1E & \seq{rcKrrrCrCCLNWR} & 1,2,4,5,6,8 & 0.679 & 4.464 & $3\text{--}<5$ \\
86 & 4K1E & \seq{rcCrGrCrQCNCDR} & 1,2,4,6,8 & 0.761 & 4.526 & $3\text{--}<5$ \\
87 & 4K1E & \seq{rQrrGrCrKCNCDR} & 1,3,4,6,8 & 1.878 & 4.594 & $3\text{--}<5$ \\
88 & 4K1E & \seq{rQCrGrCrQGFrDR} & 1,4,6,8,12 & 0.774 & 4.694 & $3\text{--}<5$ \\
89 & 4K1E & \seq{rcGrCrCrQCQCMR} & 1,2,4,6,8 & 0.763 & 4.751 & $3\text{--}<5$ \\
90 & 4K1E & \seq{rcrrGrCrFCQCGC} & 1,2,3,4,6,8 & 1.156 & 4.927 & $3\text{--}<5$ \\
91 & 4K1E & \seq{rcCrCrNrQGECDR} & 1,2,4,6,8 & 1.094 & 4.970 & $3\text{--}<5$ \\
92 & 4KEL & \seq{RcrrGNCrQGFCDR} & 2,3,4,8 & 0.654 & 2.379 & $<3$ \\
93 & 4KEL & \seq{CcCrCNQrFGQRDR} & 2,4,8 & 1.188 & 2.477 & $<3$ \\
94 & 4KEL & \seq{RcrrGrCWFCQRGR} & 2,3,4,6 & 0.988 & 3.421 & $3\text{--}<5$ \\
95 & 4KEL & \seq{RrGrrrRrQGrRDR} & 2,4,5,6,8,11 & 0.695 & 4.174 & $3\text{--}<5$ \\
96 & 4KEL & \seq{RcrrrQCrQGrCLR} & 2,3,4,5,8,11 & 0.657 & 4.267 & $3\text{--}<5$ \\
97 & 4KEL & \seq{RrrrCQCQrGLCGR} & 2,3,4,9 & 0.666 & 4.367 & $3\text{--}<5$ \\
98 & 4KEL & \seq{RcrKGrNrQGQCDR} & 2,3,6,8 & 1.426 & 4.595 & $3\text{--}<5$ \\
99 & 4KEL & \seq{RcCrGQNrQGQCGR} & 2,4,8 & 1.061 & 4.715 & $3\text{--}<5$ \\
100 & 4KEL & \seq{CGrrGrCrQGrCGR} & 3,4,6,8,11 & 1.545 & 4.871 & $3\text{--}<5$ \\
101 & 4KEL & \seq{RcrrrrCrFCrRDR} & 2,3,4,5,6,8,11 & 0.679 & 4.984 & $3\text{--}<5$ \\

\end{longtable}
\endgroup
\end{landscape}
